\documentclass[10pt]{article}
% TMLR style (https://github.com/JmlrOrg/tmlr-style-file).
% [preprint] shows the author names and drops the "Under review as submission to TMLR" header.
\usepackage[preprint]{tmlr}

\usepackage{hyperref}
\usepackage{url}
\usepackage{graphicx}
\usepackage{booktabs}
\usepackage{array}
\newcolumntype{L}[1]{>{\raggedright\arraybackslash}p{#1\linewidth}}

\title{EvoSym: Evolvable Neuro-Symbolic Intelligence}

% Separate authors with \AND.
\author{\name Gian Seifert \\
      \addr University of St.\ Gallen
      \AND
      \name Cyril Gabriele \\
      \addr University of St.\ Gallen
      \AND
      \name Filipp Rechsteiner \\
      \addr University of St.\ Gallen}

\begin{document}

\maketitle

\begin{abstract}
EvoSym answers refund questions about delayed SBB trains, asked in the passenger's own words. It couples GPT-5 mini to a FrameX knowledge base of the SBB trains of 7 November 2026 through five coupling patterns: the model turns the message into quoted slots and phrases the reply around placeholders, while the rules identify the train, measure the delay at the destination, apply the refund policy and hand every open case to a person. An independent evaluator matched all 17 policy cases, 17 of 17 live conversations and 9 of 9 dialogue and fault scenarios, and exposed a rounding convention that the policy left open. Loading one small world per decision brought a verdict from up to 166\,s down to about 13\,s; the median live reply took 16.6\,s. The results cover one operating day.
\end{abstract}

\section{Introduction}
\label{sec:introduction}

Hackathon~2 asked us to make the course's FrameX knowledge base of Swiss passenger rail hybrid~\citep{h2task}. A passenger describes a delayed journey in their own words; the system asks for what is missing, identifies the train, measures the delay at the destination from the trains that really ran on Saturday, 7 November 2026, and either answers whether a refund is due under the SBB policy or hands the case to customer service. Every station, time, train, delay and percentage in an answer must come from the knowledge base. The design had to combine at least four coupling patterns, justify them with the decision guide, and keep the five invariants of the pattern catalogue (Table~\ref{tab:invariants})~\citep{handschuh2026patterns}. Gian Seifert served as the domain expert, Filipp Rechsteiner as the developer and Cyril Gabriele as the auditor; each section reports the three views in this order.

\section{Methods}
\label{sec:methods}

% Gian's contribution: the domain expert's view. Sources are in the sibling hybridai-hackathon02 repository.
% Mandate: FlightMode Sheet 1, item 1.2; Studio slides 20.10.2026 (Language & Grounding).
% Policy: fx/47-refund-policy.fx (commit 1f32587, rounding revised in c103afd). Contract: docs/refund-chat.md.

\subsection{Language gaps and grounding}
\label{sec:usecases}

The domain expert's mandate was to decide where free language enters and leaves the system, which facts must come from the knowledge base, and what a wrong answer costs~\citep{h2task}. In FlightMode this left automatic delay compensation, with structured inputs, to rules alone. Of the five language gaps, the assistant closes questions (the message becomes slots), dialogue (it asks for missing or ambiguous details) and explanations (the verdict and its reason are worded), but not new text sources or rule maintenance: the only text is the message, and the policy is given. A wrong ``yes'' costs SBB money and a wrong ``no'' costs the passenger, so every value comes from the rules and whatever they cannot settle goes to a person~\citep{h2refundchat}.

\subsection{Refund policy as rules}
\label{sec:policy}

The domain expert encoded the policy of Sheet~2 in \texttt{47-refund-policy.fx}~\citep{h2policy}, on top of the delay that \texttt{46-refunds.fx} measures at the destination: below 60 minutes no refund, from 60 minutes 25\,\%, from 120 minutes 50\,\%, no payout below CHF~5, and a hand-off with its reason when the arrival was not measured or the day is not in the data. Three choices follow from the domain. All policy numbers live in one fact, \texttt{refundPolicy}, so every proof shows its threshold and no threshold sits in a prompt or in Python. ``Entitled but not paid out'' stays apart from ``not entitled'': CHF~4.90 due is a claim that only the minimum blocks. And the slot \texttt{dayOnRecord} is closed, so a day outside the data is handed on while the rest of the knowledge base stays open-world. Twelve tests cover the band edges (59/60 and 119/120 minutes), a missing price, the minimum, an unmeasured arrival and a day outside the data. The language contract lists what a reply must never say, such as a verdict in a hand-off; policy, tests and contract were committed one hour into the open work, so all roles built on one spec~\citep{h2refundchat}. The domain expert also built the live workflow view, which shows what each model call and rule step sent and received.

% Filipp's contribution: the developer's view. Host: src/pipeline/. Decisions: docs/adrs/0001 and 0002.

\subsection{Coupling patterns}
\label{sec:patterns}

The developer's mandate was to choose the coupling patterns and to attach GPT-5 mini to the knowledge base so that the rules keep the last word. Following the decision guide~\citep{handschuh2026patterns}, we combined five. The model must read free text, so \textbf{P3} (Model as Translator) turns each message into six closed slots (origin, destination, date, time, ticket price, train) and never writes a query. With \textbf{P1} (Model as Sensor) every slot carries the passenger's own words as a quote, accepted only if the quote occurs in a message and the value holds against the knowledge base. \textbf{P6} (Critique and Repair) returns a failed extraction or phrasing once with a fixed reason, then asks the passenger or uses a fixed sentence. Every value the passenger reads must be exact, so with \textbf{P12} (Model Selects, Rules Write) the model picks a train only by identifier from the rule-supplied list and writes its reply around named placeholders, without seeing the values. With \textbf{P7} (Division by Competence) a case counts as decided only if the policy derived a refund status; every other case goes to customer service with the model's assessment as a note. Our FlightMode diagram had P3 or P6, P12 and P7; the build uses both and adds P1, as dates, times and prices have no closed list to check against. We did not use Rules Steer the Model (P5): with four required slots, a small closed state machine in the host decides the next step.

\subsection{Where the model attaches}
\label{sec:attach}

The model attaches in one new package, \texttt{src/pipeline/}, beside the template's layers. The rule base changed in three places: the domain expert's policy file, two edits to \texttt{46-refunds.fx} and a small gate for the hand-off. The host is the only caller of both GPT-5 mini and FrameX: it asks the model for a JSON object through the official OpenAI client and records every answer with its date, so an identical request is replayed. The model never calls the rules and never sees a query. Unlike the catalogue, where the rules issue the requests, only the assessment of an open delay check is requested by a rule, through FrameX's \texttt{competence.fx}, and carried out by the host. Table~\ref{tab:boundary} lists what crosses the boundary; the \texttt{RefundCase} and \texttt{DelayCheck} facts are written from checked identifiers only, and the host inserts the values FrameX returns into the reply word for word. A failed extraction or selection call ends the turn as a technical error without a verdict; a failed phrasing call falls back to the fixed sentence, so the verdict still reaches the passenger.

\begin{table}[t]
\caption{What crosses the model boundary, per kind of call from the host to GPT-5 mini.}
\label{tab:boundary}
\centering
\footnotesize
\setlength{\tabcolsep}{4pt}\renewcommand{\arraystretch}{0.92}
\begin{tabular}{L{0.16}L{0.30}L{0.18}L{0.29}}
\toprule
Call & To the model & From the model & Accepted only if \\
\midrule
Extract slots (P3, P1) & passenger messages, slots already known & slots with quotes & quote found and value valid; one repair (P6) \\
Ask for missing slots (P3) & conversation, names of the known values, what is missing & one question with placeholders & no digit, station name or promise of a refund \\
Select train (P12) & passenger's reply, offered trains under identifiers & one identifier or none & the identifier is in the list \\
Phrase outcome (P12) & decision, sentence with placeholders and their meaning & decision and reply & same decision, every placeholder, no digit or station name; one repair (P6) \\
Assess open case (P7) & case file, the rules' reason & verdict and reason & never: it is a note for the officer \\
\bottomrule
\end{tabular}
\end{table}

% Cyril's contribution. Sources are in the sibling hybridai-hackathon02 repository.
% Mandate: docs/cyril/mandate.md.
% Methods: docs/evaluation.md; src/eval/; tests/eval/;
% docs/adrs/0003_refund_evaluation.md and 0004_conversation_evaluation.md.
% Historical finding: docs/cyril/evaluation-baseline.md.
% Final results: docs/cyril/evaluation-e2e.md;
% config/eval/verification.yaml; eval-results/e2e-snapshot/summary.md;
% eval-results/e2e-dialogue/summary.md; eval-results/policy-integrated/summary.md.


\subsection{Independent evaluation}

The auditor's mandate was to check that the model cannot execute unchecked rules or queries, that acceptance depends on evidence, and that the design keeps the five invariants. Following Chapter~9 of \citet{albada2025agents}, the auditor evaluated rule correctness (component) separately from complete conversations (holistic). Expected answers stay outside the system under test: a policy oracle calculates compensation in Python, and a journey oracle enumerates direct and one-change routes from the frozen SBB facts. FrameX results are compared with these expectations and never generate them; snapshot hashes and route witnesses preserve the evidence. The conversation adapter invokes the same assistant as the user interface, without exposing expected answers or oracle-selected trains~\citep{h2evaluation}.

\subsection{Checks at the model boundaries}

Boundaries are checked through structured events emitted where production validation, execution and acceptance occur. An execution passes only if an earlier approval matches its operation and arguments. Accepted values must reference recorded evidence and accepted quotes must occur in the message, which establishes traceability but not that the model interpreted the words correctly. Dialogue cases cover missing and ambiguous stations, an unsupported date, rule injection, fabricated quotes, a corrected destination, exhausted repair and model or tool timeouts, and the one-repair budgets. Errors must remain errors, clarification and hand-off carry no verdict, and missing evidence stays unverified rather than counting as proof of safe execution. Table~\ref{tab:invariants} maps each invariant to its risk, protection and evidence. Latency is measured with a monotonic clock, excluding recorded or injected responses, and a separate GPT-5 mini judge scores five reply-quality criteria apart from correctness~\citep{h2evaluation}.

\begin{table}[t]
\caption{Where the design endangers each invariant, and how it is protected and checked.}
\label{tab:invariants}
\centering
\footnotesize
\setlength{\tabcolsep}{4pt}\renewcommand{\arraystretch}{0.92}
\begin{tabular}{L{0.18}L{0.23}L{0.31}L{0.21}}
\toprule
Invariant & Endangered by & Protection & Evidence \\
\midrule
1 Model output is data & a message that injects a rule or query & closed slots; facts written only from checked identifiers; execution needs a matching approval & rule-injection case \\
2 Dated observation & an unrecorded model call & every answer stored with its date and replayed & reports keep model, timestamps, transcripts \\
3 Rules decide & a value computed by the host or the model & values only from rule bindings, inserted by the host (P12) & 17/17 policy, 7/7 workspace checks \\
4 Evidence governs acceptance & a fabricated quote & quote must occur in the message; accepted values reference evidence & fabricated-quote cases \\
5 Failures stay distinct & a timeout read as ``no refund'' & errors, clarification and hand-off kept apart; no verdict on failure & 2 expected technical errors \\
\bottomrule
\end{tabular}
\end{table}


\section{Results}
\label{sec:results}

% Gian's results. Counts recomputed on 16 November 2026 from fx/70-ist-daten-sbb.fx.

\subsection{The policy against tests and data}
\label{sec:datacheck}

All 19 tests in \texttt{tests/test\_refunds.py} pass. On 7 November the largest arrival delay was 71 minutes; of 56,675 recorded arrivals, two were 60 or more minutes late and none 120 or more, so a 50\,\% refund arises only through a cancellation or a missed connection. The task's own example matches four departures within 30 minutes, all checked ones on time: the assistant must ask which train, and the answer is no refund.

% Filipp's results. Source: docs/adrs/0001_refund_case_worlds.md.

\subsection{One world per decision}

Adding a case to one session of the full workspace made the engine recompute everything: 27\,s for the first case and 166\,s by the third. Each decision therefore loads a fresh engine with the unchanged rules and only the trips they can use for that journey; listing the trains then takes about 5\,s and a verdict about 13\,s, against about 3\,s for a phrased reply, so the cost lies in the knowledge base, not the model~\citep{rechsteiner2026adr0001,rechsteiner2026adr0002}. The reduced world matched a world with every trip at the origin or destination: the same train lists on four routes, the same result in nine of nine delay checks. It also exposed a data fault: two services under one trip reference formed a Z\"urich HB--Bern train that arrived before it departed, hid every real departure and earned a false 25\,\% refund, until the rule required arrival after departure.

% Cyril's results. Sources: docs/cyril/evaluation-baseline.md; docs/cyril/evaluation-e2e.md.

\subsection{Audit findings, correctness and latency}
\label{sec:audit-results}

The baseline exposed four rounding mismatches among 17 policy cases: production discarded fractional Rappen, while the agreed oracle compared the exact amount with the CHF~5 minimum and rounded only the displayed amount half-up, so 25\,\% of CHF~19.99 shows as CHF~5.00 but must not be paid. The final implementation matched all 17 policy cases~\citep{h2baseline,h2evaluation}. The recorded evaluation of 10 November reports 17 of 17 correct live conversations, all passing every instrumented invariant, and 9 of 9 dialogue and fault scenarios, two of them expected technical errors. Seven cases, including cancellation and missing arrival, also matched on the full FrameX workspace. The median live response took 16.6\,s and the 95th percentile 25.2\,s over 17 sequential requests, without judge time, and all 24 eligible replies were scored~\citep{h2verification}.


\section{Discussion}
\label{sec:discussion}

% Gian's contribution to the Discussion: design choices in hindsight, from the domain expert's view.
% Sources: hackathon-day session notes (gaps spotted around 16:15); docs/cyril/evaluation-baseline.md (rounding);
% data counts recomputed on 16 November 2026 from fx/70-ist-daten-sbb.fx.

\subsection{Design choices in hindsight}

Four domain choices would be made differently. First, rounding: the policy sheet does not say how to round fractional Rappen. The domain expert chose floor rounding, the auditor's oracle half-up, and the evaluation found four mismatches (Section~\ref{sec:results}). A convention the sheet leaves open belongs in the spec, agreed by all roles before code is written. Second, relative dates such as ``last Saturday'' and the refundable share of a return ticket were noticed as gaps but never specified, so nobody decided them. Third, the ticket price is the only input the knowledge base cannot confirm: the quote check proves that the passenger wrote it, not that they paid it. Answering with the percentage only and leaving the amount to SBB, which holds the ticket record, would remove that input. Fourth, refund cases are rare on the operating day; choosing the real refund journeys in FlightMode would have given all roles the same test cases from the start.

One choice held up. In the FlightMode debate the developer argued that rules steering the model (P5) would need fewer model calls, but the measured cost lay in the knowledge base (about 13\,s per verdict), not the model (about 3\,s per reply)~\citep{h2adrs}.

% Filipp's discussion: limits of the coupling and lessons learned.

\subsection{Limits of the coupling}

The next step is decided by host code, not by a FrameX flow, so the step order and the host's checks appear in the step trace but not in the proofs, and so are hand-offs raised before a delay check exists, such as an unidentified train. P12 protects values, not wording: a fluent but misleading sentence can pass, so the host adds its own line with the checked times below each verdict. The trip selection repeats the rules' search bounds in Python and was compared on a few routes only. Each world needs about 2\,GB of memory, so decisions run one at a time. \emph{Lessons learned:} data faults surface as plausible verdicts that no check on the model side catches, so the rules should define a valid trip; and once a world per decision is cheap enough, a FrameX flow (P5) would put the step order into the proofs~\citep{rechsteiner2026adr0001,rechsteiner2026adr0002}.

% Cyril's discussion: evaluation limits and lessons learned.

\subsection{Evaluation limits}

Latency comes from a small sequential sample, not concurrent load; the judge scores are uncalibrated proxies, not deterministic checks; and the audit covers one operating day and journeys with at most one change. It cannot establish source-data accuracy, user satisfaction or repeated-run consistency~\citep{h2verification,h2evaluation}. \emph{Lessons learned:} only an oracle outside the system exposed the rounding gap, which also avoids Circular Evaluation~\citep{handschuh2026patterns}; audit evidence belongs where the checks run, since parsing the workflow view could invent approvals~\citep{h2adr0004}; and a judge other than GPT-5 mini, or a rule-based read-back, would remove the Self-Judge risk.


\bibliography{references}
\bibliographystyle{tmlr}

\appendix
\section{Appendix}

\subsection{Use of coding agents}

Coding agents were used only after FlightMode. Gian used Claude Code with Claude Opus 5.5 for coding during the hackathon and for report writing.

Filipp used Claude Code with Claude Opus 5.5 for coding during the hackathon, and with Claude Fable 5.2 to draft the developer's sections of this report from the repository.

Cyril used Codex with GPT-6 Sol for coding during the hackathon and for report writing.

\subsection{Declaration of authorship}

We declare that we wrote this report ourselves, using only the sources and aids listed. Contributions of coding agents are declared above.

\end{document}
